% ================================================================
%  SESIÓN 13 — De tabla a gráfica
%  Almería en Cifras y Formas · Matemáticas 1.º ESO
%  IES Al-Ándalus · 2025-2026
% ================================================================
\documentclass[aspectratio=169, 12pt, spanish]{beamer}
\usepackage{almeria-beamer}

\renewcommand{\SessionNumber}{13}
\renewcommand{\SessionTitle}{De tabla a gráfica}
\renewcommand{\SessionName}{Sesión 13}
\renewcommand{\SessionObjective}{Representar datos de una tabla en una gráfica (diagrama de barras y gráfica lineal) eligiendo el tipo adecuado según el tipo de variable}
\renewcommand{\BlockName}{Bloque 3: Datos y gráficas}

\begin{document}

% ---------------------------------------------------------------
%  PORTADA
% ---------------------------------------------------------------
\begin{frame}[plain]
  \titlepage
\end{frame}

% ---------------------------------------------------------------
%  FRAME 1 — ¿Por qué representar gráficamente?
% ---------------------------------------------------------------
\begin{frame}{¿Por qué usar gráficas?}
  \begin{columns}[T]
    \begin{column}{0.52\textwidth}
      \begin{teoriabox}
        \small
        Una tabla muestra los datos con \textbf{precisión},
        pero una gráfica los hace \textbf{visibles}:
        \begin{itemize}
          \item Revela \textbf{tendencias} (subida, bajada, estabilidad).
          \item Permite detectar \textbf{valores atípicos} de un vistazo.
          \item Facilita la \textbf{comparación} entre series.
          \item Comunica la información más rápido que un texto.
        \end{itemize}
      \end{teoriabox}
    \end{column}
    \begin{column}{0.44\textwidth}
      \begin{notabox}[Datos → Gráfica → Conclusión]
        \small
        \begin{center}
        \begin{tikzpicture}[node distance=0.7cm,
            caja/.style={rounded corners=4pt, draw=AlmAzulM,
                         fill=AlmFondo, text width=2.2cm,
                         align=center, font=\tiny\bfseries}]
          \node[caja] (d) {Tabla\\de datos};
          \node[caja, below=of d] (g) {Gráfica\\adecuada};
          \node[caja, below=of g] (c) {Conclusión\\estadística};
          \draw[-Stealth, AlmAzulM] (d)--(g);
          \draw[-Stealth, AlmAzulM] (g)--(c);
        \end{tikzpicture}
        \end{center}
      \end{notabox}
    \end{column}
  \end{columns}
\end{frame}

% ---------------------------------------------------------------
%  FRAME 2 — Tipos de gráficas y cuándo usarlas
% ---------------------------------------------------------------
\begin{frame}{Tipos de gráficas y cuándo usarlas}
  \small
  \renewcommand{\arraystretch}{1.4}
  \begin{tabular}{>{\bfseries\color{AlmAzulM}}l l l}
    \toprule
    \rowcolor{AlmAzul!12}
    \textbf{Tipo de gráfica} & \textbf{Cuándo usarla} & \textbf{Ejemplo en Almería} \\
    \midrule
    Diagrama de barras
      & Variable $x$ discreta o categórica
      & Visitantes por barrio \\
    \rowcolor{AlmFondo}
    Gráfica lineal
      & $x$ continua o temporal; muestra tendencia
      & Temperatura media mensual \\
    Diagrama de puntos
      & Dos variables cuantitativas; correlación
      & Temperatura vs.\ turistas \\
    \rowcolor{AlmFondo}
    Diagrama de sectores
      & Partes de un todo (porcentajes)
      & Tipo de transporte usado \\
    \bottomrule
  \end{tabular}

  \vspace{0.3cm}
  \begin{notabox}[Regla práctica]
    \small
    Si $x$ son \textbf{categorías o meses} que comparo → barras.\\
    Si $x$ es \textbf{tiempo continuo} y quiero ver la evolución → línea.
  \end{notabox}
\end{frame}

% ---------------------------------------------------------------
%  FRAME 3 — Elementos de una gráfica (TikZ anotado)
% ---------------------------------------------------------------
\begin{frame}{Elementos de una gráfica bien construida}
  \begin{columns}[T]
    \begin{column}{0.54\textwidth}
      \begin{center}
      \begin{tikzpicture}[font=\scriptsize, scale=0.88]
        % Ejes
        \draw[AlmAzulM, line width=1.2pt, -Stealth] (0,0)--(6.2,0)
          node[right, color=AlmAzulM, font=\tiny\bfseries] {Mes ($x$)};
        \draw[AlmAzulM, line width=1.2pt, -Stealth] (0,0)--(0,4.2)
          node[above, color=AlmAzulM, font=\tiny\bfseries,
               xshift=-0.3cm, rotate=90] {Visitantes ($y$)};
        % Barras de muestra
        \foreach \x/\h/\c in {
          0.6/0.9/AlmAzulC, 1.2/1.0/AlmAzulC, 1.8/1.3/AlmAzulC,
          2.4/1.9/AlmAzulM, 3.0/2.7/AlmAzulM}{
          \fill[\c, opacity=0.7] (\x,0) rectangle (\x+0.5,\h);
        }
        % Título
        \node[color=AlmAzul, font=\tiny\bfseries] at (3.0,4.0)
          {Visitantes mensuales a la Alcazaba};
        % Escala y
        \foreach \h/\v in {1.0/10000, 2.0/20000, 3.0/30000}{
          \draw[AlmGris!40, line width=0.4pt] (0,\h)--(5.8,\h);
          \node[left, color=AlmGris, font=\tiny] at (0,\h) {\v};
        }
        % Etiquetas x
        \foreach \x/\l in {0.85/E, 1.45/F, 2.05/M, 2.65/A, 3.25/M}{
          \node[below=2pt, color=AlmGris] at (\x,0) {\l};
        }
        % Anotaciones externas
        \draw[-Stealth, AlmTerra, thin] (5.4,3.8) -- (4.6,2.8)
          node[above right, color=AlmTerra, font=\tiny\bfseries,
               xshift=8pt, yshift=4pt] {Título};
        \draw[-Stealth, AlmVerde, thin] (5.4,0.5) -- (0.7,0.5)
          node[above right, color=AlmVerde, font=\tiny\bfseries,
               xshift=-24pt, yshift=6pt] {Barra};
        \node[AlmAzulM, font=\tiny\bfseries] at (-1.1,2.0)
          {\rotatebox{90}{Escala $y$}};
        \node[AlmAzulM, font=\tiny\bfseries] at (3.2,-0.5)
          {Etiquetas $x$};
      \end{tikzpicture}
      \end{center}
    \end{column}
    \begin{column}{0.42\textwidth}
      \begin{teoriabox}
        \small
        Todo gráfico debe tener:
        \begin{itemize}
          \item \textbf{Eje $x$}: variable independiente + unidad
          \item \textbf{Eje $y$}: variable dependiente + unidad
          \item \textbf{Escala}: intervalos regulares, comenzando en~0
          \item \textbf{Título}: qué se representa
          \item \textbf{Leyenda}: si hay más de una serie
        \end{itemize}
      \end{teoriabox}
    \end{column}
  \end{columns}
\end{frame}

% ---------------------------------------------------------------
%  FRAME 4 — Pasos para dibujar un gráfico
% ---------------------------------------------------------------
\begin{frame}{Pasos para trazar una gráfica}
  \begin{columns}[T]
    \begin{column}{0.48\textwidth}
      \begin{pasobox}[1]
        \small Decide el \textbf{tipo} de gráfica
        (¿barras o línea?).
      \end{pasobox}\vspace{3pt}
      \begin{pasobox}[2]
        \small Dibuja los \textbf{ejes} $x$ (horizontal) e $y$ (vertical).
      \end{pasobox}\vspace{3pt}
      \begin{pasobox}[3]
        \small \textbf{Etiqueta} cada eje: variable + unidad.
      \end{pasobox}\vspace{3pt}
      \begin{pasobox}[4]
        \small Elige la \textbf{escala}: mayor valor $\times\;1{,}1$
        para el máximo del eje.
      \end{pasobox}
    \end{column}
    \begin{column}{0.48\textwidth}
      \begin{pasobox}[5]
        \small \textbf{Representa} cada par $(x,y)$
        de la tabla como punto.
      \end{pasobox}\vspace{3pt}
      \begin{pasobox}[6]
        \small \textbf{Une} los puntos con línea (gráfica lineal)
        o rellena barras (diagrama de barras).
      \end{pasobox}\vspace{3pt}
      \begin{pasobox}[7]
        \small Añade el \textbf{título} descriptivo.
      \end{pasobox}

      \vspace{0.25cm}
      \begin{notabox}[Clave del paso 4]
        \small
        El eje $y$ \textbf{siempre empieza en 0}
        salvo razón justificada.
      \end{notabox}
    \end{column}
  \end{columns}
\end{frame}

% ---------------------------------------------------------------
%  FRAME 5 — Diagrama de barras: visitantes Almería ene–jun
% ---------------------------------------------------------------
\begin{frame}{Diagrama de barras: turistas en Almería (ene–jun)}
  \begin{center}
  \begin{tikzpicture}
    \begin{axis}[
      almeria bar,
      width=0.82\linewidth, height=0.56\linewidth,
      title={Turistas mensuales en Almería (enero–junio)},
      xlabel={Mes},
      ylabel={Turistas},
      symbolic x coords={Enero,Febrero,Marzo,Abril,Mayo,Junio},
      xtick=data,
      xticklabel style={font=\tiny, rotate=15, anchor=north east},
      ymin=0, ymax=330000,
      ytick={0,50000,100000,150000,200000,250000,300000},
      yticklabels={0,50\,k,100\,k,150\,k,200\,k,250\,k,300\,k},
      nodes near coords style={font=\tiny\bfseries, color=AlmAzul,
                                rotate=90, anchor=west},
    ]
      \addplot[fill=AlmAzulM!75, draw=AlmAzul, fill opacity=0.85] coordinates {
        (Enero,60000)
        (Febrero,65000)
        (Marzo,85000)
        (Abril,130000)
        (Mayo,180000)
        (Junio,280000)
      };
    \end{axis}
  \end{tikzpicture}
  \end{center}
  {\small\color{AlmGris}
   Usamos barras porque los meses son \textbf{categorías discretas}
   que queremos \textit{comparar}.}
\end{frame}

% ---------------------------------------------------------------
%  FRAME 6 — Gráfica lineal: temperatura mensual en Almería
% ---------------------------------------------------------------
\begin{frame}{Gráfica lineal: temperatura media mensual en Almería}
  \begin{center}
  \begin{tikzpicture}
    \begin{axis}[
      almeria line,
      width=0.82\linewidth, height=0.54\linewidth,
      title={Temperatura media mensual en Almería (°C)},
      xlabel={Mes},
      ylabel={Temperatura (°C)},
      xtick={1,2,3,4,5,6,7,8,9,10,11,12},
      xticklabels={E,F,M,A,M,J,J,A,S,O,N,D},
      ymin=8, ymax=34,
      ytick={10,15,20,25,30},
    ]
      \addplot[
        color=AlmTerra, line width=2pt, mark=*, mark size=2.5pt,
        mark options={fill=AlmTerra!70, draw=AlmTerra}
      ] coordinates {
        (1,12)(2,13)(3,16)(4,19)(5,23)(6,27)
        (7,30)(8,30)(9,26)(10,21)(11,16)(12,13)
      };
      \node[font=\tiny\bfseries, color=AlmTerra, above=4pt] at (axis cs:7,30)
        {30\,°C};
      \node[font=\tiny\bfseries, color=AlmAzulM, below=4pt] at (axis cs:1,12)
        {12\,°C};
    \end{axis}
  \end{tikzpicture}
  \end{center}
  {\small\color{AlmGris}
   Usamos línea porque la temperatura varía de forma \textbf{continua}
   a lo largo del tiempo y queremos ver la \textit{tendencia}.}
\end{frame}

% ---------------------------------------------------------------
%  FRAME 7 — Barras vs. línea: ¿cuándo usar cada una?
% ---------------------------------------------------------------
\begin{frame}{Barras vs. línea: la elección correcta}
  \begin{columns}[T]
    \begin{column}{0.48\textwidth}
      \begin{definicionbox}{Diagrama de barras — úsalo cuando:}
        \small
        \begin{itemize}
          \item $x$ son \textbf{categorías separadas}\\
                (tipos, barrios, meses como etiquetas)
          \item Quieres \textbf{comparar magnitudes}
          \item No hay continuidad entre valores de $x$
        \end{itemize}
        \vspace{4pt}
        \textit{Ejemplo}: nº de plazas hoteleras por barrio de Almería
      \end{definicionbox}
    \end{column}
    \begin{column}{0.48\textwidth}
      \begin{definicionbox}{Gráfica lineal — úsala cuando:}
        \small
        \begin{itemize}
          \item $x$ es \textbf{continua o temporal}
                (tiempo, distancia, temperatura)
          \item Quieres mostrar una \textbf{tendencia} o evolución
          \item Tiene sentido hablar de valores intermedios
        \end{itemize}
        \vspace{4pt}
        \textit{Ejemplo}: temperatura media horaria en un día de verano
      \end{definicionbox}
    \end{column}
  \end{columns}

  \vspace{0.3cm}
  \begin{notabox}[¿Y si tengo dudas?]
    \small
    Pregúntate: «¿Tendría sentido un valor \textit{entre} dos puntos del eje $x$?»
    Si sí → línea. Si no → barras.
  \end{notabox}
\end{frame}

% ---------------------------------------------------------------
%  FRAME 8 — Ejercicio RECORDAR
% ---------------------------------------------------------------
\begin{frame}{Ejercicio: tipos de gráficas}
  \begin{recordarbox}
    \begin{enumerate}
      \item Nombra \textbf{4 tipos de gráficas} que conoces.
      \item Para cada una, escribe \textbf{un ejemplo concreto}
            relacionado con Almería en el que sea apropiada.
      \item Completa la tabla:
    \end{enumerate}

    \vspace{0.2cm}
    \small
    \renewcommand{\arraystretch}{1.3}
    \begin{tabular}{lcc}
      \toprule
      \rowcolor{AlmAzul!12}
      \textbf{Situación} & \textbf{¿Barras o línea?} & \textbf{¿Por qué?} \\
      \midrule
      Visitantes por mes (etiquetas) & & \\
      \rowcolor{AlmFondo}
      Temperatura a lo largo del día & & \\
      Cuota de turistas por país & & \\
      \rowcolor{AlmFondo}
      Evolución anual de pernoctaciones & & \\
      \bottomrule
    \end{tabular}
  \end{recordarbox}
\end{frame}

% ---------------------------------------------------------------
%  FRAME 9 — Ejercicio COMPRENDER
% ---------------------------------------------------------------
\begin{frame}{Ejercicio: ¿barras o línea?}
  \begin{comprenderbox}
    Se recogen estos datos sobre alquileres en Almería:

    \vspace{0.2cm}
    \begin{center}
    \small
    \begin{tabular}{cc}
      \toprule
      \rowcolor{AlmAzul!12}
      \textbf{Distancia al centro (km)} & \textbf{Precio alquiler (€/m²)} \\
      \midrule
      0{,}5 & 12 \\
      \rowcolor{AlmFondo}
      1{,}0 & 10 \\
      1{,}5 & 8{,}5 \\
      \rowcolor{AlmFondo}
      2{,}0 & 7 \\
      2{,}5 & 6 \\
      \bottomrule
    \end{tabular}
    \end{center}

    \vspace{0.2cm}
    \begin{enumerate}
      \item ¿Deberías usar un diagrama de barras o una gráfica lineal?
            \textbf{Justifica tu respuesta}.
      \item ¿Cuál es la variable independiente? ¿Y la dependiente?
      \item Describe en una frase la tendencia que observas.
    \end{enumerate}
  \end{comprenderbox}
\end{frame}

% ---------------------------------------------------------------
%  FRAME 10 — Ejercicio APLICAR
% ---------------------------------------------------------------
\begin{frame}{Ejercicio: graficar los datos del autobús}
  \begin{aplicarbox}
    Retoma la tabla del autobús de Almería (sesión 12):

    \vspace{0.2cm}
    \begin{center}
    \small
    \begin{tabular}{cc}
      \toprule
      \rowcolor{AlmAzul!12}
      \textbf{Paradas (\(x\))} & \textbf{Tiempo min (\(y\))} \\
      \midrule
      2 & 5 \\ \rowcolor{AlmFondo}
      4 & 10 \\
      6 & 15 \\ \rowcolor{AlmFondo}
      8 & 20 \\
      12 & 30 \\
      \bottomrule
    \end{tabular}
    \end{center}

    \vspace{0.2cm}
    \begin{enumerate}
      \item Dibuja la gráfica lineal correspondiente en papel cuadriculado.
            Etiqueta ambos ejes con nombre y unidad.
            Elige una escala adecuada para el eje $y$ (máx. 30~min).
      \item ¿Qué \textbf{tendencia} observas?
      \item ¿Qué representa la \textbf{inclinación} de la línea?
    \end{enumerate}
  \end{aplicarbox}
\end{frame}

% ---------------------------------------------------------------
%  FRAME 11 — Error frecuente
% ---------------------------------------------------------------
\begin{frame}{¡Cuidado con este error!}
  \begin{errorbox}
    \textbf{Empezar el eje \(y\) en un valor distinto de cero}
    para exagerar las diferencias:

    \vspace{0.3cm}
    \begin{columns}[T]
      \begin{column}{0.46\textwidth}
        \incorrecto\; \textbf{Eje \(y\) desde 95\,000}\\[4pt]
        \begin{center}
        \begin{tikzpicture}[scale=0.7]
          \begin{axis}[
            width=0.9\linewidth, height=0.7\linewidth,
            axis lines=left, ymin=95000, ymax=105000,
            xtick={1,2,3}, xticklabels={E,F,M},
            ytick={95000,100000,105000},
            yticklabels={95k,100k,105k},
            tick label style={font=\tiny},
            label style={font=\tiny},
            ylabel style={font=\tiny},
            ylabel={Visitantes},
            xlabel={Mes},
            ybar, bar width=0.4cm,
            nodes near coords, nodes near coords style={font=\tiny},
          ]
            \addplot[fill=AlmTerra!60] coordinates
              {(1,97000)(2,100000)(3,104000)};
          \end{axis}
        \end{tikzpicture}
        \end{center}
        \small\color{AlmTerra}
        ¡Parece que la diferencia es enorme!
      \end{column}
      \begin{column}{0.46\textwidth}
        \correcto\; \textbf{Eje \(y\) desde 0}\\[4pt]
        \begin{center}
        \begin{tikzpicture}[scale=0.7]
          \begin{axis}[
            width=0.9\linewidth, height=0.7\linewidth,
            axis lines=left, ymin=0, ymax=115000,
            xtick={1,2,3}, xticklabels={E,F,M},
            ytick={0,50000,100000},
            yticklabels={0,50k,100k},
            tick label style={font=\tiny},
            label style={font=\tiny},
            ylabel style={font=\tiny},
            ylabel={Visitantes},
            xlabel={Mes},
            ybar, bar width=0.4cm,
            nodes near coords, nodes near coords style={font=\tiny},
          ]
            \addplot[fill=AlmVerde!60] coordinates
              {(1,97000)(2,100000)(3,104000)};
          \end{axis}
        \end{tikzpicture}
        \end{center}
        \small\color{AlmVerde}
        La diferencia es real pero pequeña.
      \end{column}
    \end{columns}
  \end{errorbox}
\end{frame}

% ---------------------------------------------------------------
%  FRAME 12 — Evidencia del día
% ---------------------------------------------------------------
\begin{frame}{Evidencia del día}
  \begin{evidenciabox}
    Dibuja en tu cuaderno \textbf{dos gráficas} a partir de las tablas
    que construiste en la sesión anterior:

    \vspace{0.15cm}
    \begin{enumerate}
      \item \textbf{Gráfica A} — Temperaturas medias de Almería:
            \begin{itemize}
              \small
              \item Elige el tipo correcto (barras o línea).
              \item Etiqueta los ejes: «Mes» y «Temperatura (°C)».
              \item Escala: 0–35\,°C.
              \item Título: «Temperatura media mensual en Almería».
            \end{itemize}
      \item \textbf{Gráfica B} — Precio del taxi en Almería:
            \begin{itemize}
              \small
              \item Elige el tipo correcto.
              \item Etiqueta: «Km recorridos» y «Precio (€)».
              \item Escala acorde al máximo de la tabla.
            \end{itemize}
    \end{enumerate}

    \vspace{0.15cm}
    \small
    Para cada gráfica: justifica en una frase por qué elegiste
    ese tipo de representación.
  \end{evidenciabox}
\end{frame}

% ---------------------------------------------------------------
%  FRAME 13 — Resumen de sesión
% ---------------------------------------------------------------
\begin{frame}{Resumen: lo que hemos aprendido hoy}
  \begin{resumenbox}
    \begin{columns}[T]
      \begin{column}{0.48\textwidth}
        \textbf{Elegir el tipo de gráfica:}
        \begin{itemize}
          \item \textbf{Barras}: categorías, comparar
          \item \textbf{Línea}: variable continua, tendencia
          \item \textbf{Puntos}: correlación entre dos variables
          \item \textbf{Sectores}: partes de un todo
        \end{itemize}
      \end{column}
      \begin{column}{0.48\textwidth}
        \textbf{Elementos imprescindibles:}
        \begin{itemize}
          \item Ejes etiquetados con unidades
          \item Escala regular comenzando en 0
          \item Título claro y descriptivo
          \item Leyenda si hay varias series
        \end{itemize}
      \end{column}
    \end{columns}

    \vspace{0.25cm}
    \separador
    \vspace{0.15cm}
    \small
    \textbf{Error crítico:} empezar el eje $y$ en un valor
    distinto de cero distorsiona visualmente las diferencias y
    puede ser \textit{engañoso}.
  \end{resumenbox}

  \vspace{0.2cm}
  \begin{center}
    \bloomtag{Recordar}\quad
    \bloomtag{Comprender}\quad
    \bloomtag{Aplicar}
  \end{center}
\end{frame}

\end{document}
